\section{Intégration Système : Processeur et Mémoire}

\subsection{Modification du processeur}

Le processeur a été modifié pour déporter son stockage de données vers la mémoire externe \texttt{quadram}, en intégrant les bus \texttt{mem\_addr}, \texttt{mem\_wdata} et \texttt{mem\_rdata}, ainsi que les signaux de contrôle \texttt{mem\_we}, \texttt{clk} et \texttt{rst}.

\begin{lstlisting}[caption={Interface externe du processeur}, label={lst:processor_ram}]
module system (
  input wire clk, rst,
  output reg  [11:0] mem_addr,
  output reg  [31:0] mem_wdata,
  input  wire [31:0] mem_rdata,
  output reg         mem_we
);
\end{lstlisting}

\subsection{Création d'un bloc Top-Level}

Un composant \textit{Top-Level} rédigé en VHDL relie le processeur (en Verilog) et le bloc \texttt{QuadRAM}, en assurant la conversion des types entre les signaux standards et les vecteurs non signés.

\begin{lstlisting}[caption={Extrait du Top-Level système}, label={lst:top_level}]
entity system_top is
    Port (
        clk : in STD_LOGIC;
        rst : in STD_LOGIC
    );
end entity system_top;

architecture RTL of system_top is

    signal mem_addr  : std_logic_vector(11 downto 0);
    signal mem_wdata : std_logic_vector(31 downto 0);
    signal mem_rdata : std_logic_vector(31 downto 0);
    signal mem_we    : std_logic;

    signal quadram_addr  : unsigned(11 downto 0);
    signal quadram_wdata : unsigned(31 downto 0);
    signal quadram_rdata : unsigned(31 downto 0);

begin

    quadram_addr  <= unsigned(mem_addr);
    quadram_wdata <= unsigned(mem_wdata);
    mem_rdata     <= std_logic_vector(quadram_rdata);

    processor_inst : entity work.system
        port map (
            clk       => clk,
            rst       => rst,
            mem_addr  => mem_addr,
            mem_wdata => mem_wdata,
            mem_rdata => mem_rdata,
            mem_we    => mem_we
        );

    quadram_inst : entity work.quadram
        generic map (
            ADDR_SELECT => 2,
            ADDR_LENGHT => 12,
            DATA_LENGHT => 32
        )
        port map (
            clock    => clk,
            we       => mem_we,
            address  => quadram_addr,
            datain   => quadram_wdata,
            dataout  => quadram_rdata
        );

end architecture RTL;
\end{lstlisting}

\subsection{Validation du système intégral}

Dans le même esprit que pour la validation du modèle processeur, nous validons notre système intégral en simulant l'exécution d'un programme de test. 
Le chronogramme ci-dessous (Figure~\ref{fig:processor_ram}) illustre le fonctionnement du processeur en interaction avec la mémoire \texttt{quadram}.
\\

\begin{itemize}
    \item \textcolor{red}{\textbf{Étape 1}} : le processeur additionne \texttt{x1} = 2 et \texttt{x2} = 4, puis écrit le résultat 6 à l'adresse mémoire 0.
    \item \textcolor{blue}{\textbf{Étape 2}} : il multiplie \texttt{x1} = 2 par \texttt{x2} = 4, puis écrit le résultat 8 à l'adresse mémoire 1.
    \item \textcolor{magenta}{\textbf{Étape 3}} : il divise \texttt{x2} = 4 par \texttt{x1} = 2, puis écrit le résultat 2 à l'adresse mémoire 2.
    \item \textcolor{orange}{\textbf{Étape 4}} : il lit le contenu de l'adresse mémoire 0 et charge la valeur 6 dans \texttt{x2}.
\end{itemize}

Voici le chronogramme obtenu (Figure~\ref{fig:processor_ram}) :

\begin{figure}[H]
    \centering
    \includegraphics[width=0.99\textwidth]{4_OPERATION_modified.png}
    \caption{Opérations arithmétiques et échanges mémoire du processeur avec la \texttt{quadram}}
    \label{fig:processor_ram}
\end{figure}

Les repères colorés du chronogramme suivent ces quatre étapes : le rouge identifie l'addition et l'écriture de 6 à l'adresse 0 ; le bleu, la multiplication et l'écriture de 8 à l'adresse 1 ; le rose, la division et l'écriture de 2 à l'adresse 2 ; enfin, l'orange correspond à la lecture de l'adresse 0 et au chargement de la valeur 6 dans \texttt{x2}.